% Integration draft for the frozen selection-v2 and matched-v2 studies.
% Result-bearing tokens must be replaced only from completed, verified artifacts.
% This file deliberately contains no formal Q2 or Q3--Q4 result values.

% BEGIN SELECTION BLOCK selection_protocol
\subsection{Validation-selected predictive comparison}\label{sec:selection}
{\color{blue}We compare TL-QLDMR, its no-variance and no-discrepancy variants, pooled weighted support vector quantile regression (SVQR), target-only SVQR, and histogram gradient boosting with quantile loss (HistGBR). The six within-site normal-to-extreme tasks form the primary scope, and the six leave-one-farm-out tasks form a secondary transfer-stress scope. Every site and method is retained in its prespecified scope. Within each primary task, the source observations are normal-domain windows from the target farm; within each secondary task, normal-domain source observations come from the other five farms.}

{\color{blue}Each task uses exactly 128 extreme-domain target fitting labels selected deterministically from 2019. The cells are training month crossed with mutually exclusive cold, hot, and high-wind/cut-out categories, assigned in that priority order from target-time meteorological covariates within the fitting period only; power values are not used in acquisition. Each nonempty cell first receives one observation; the remaining labels are allocated by the largest-remainder rule in proportion to the cells' remaining capacities, and observations within a cell are chosen near equally spaced clock-time targets. A primary task uses at most 7,500 normal-domain source windows from the same farm, deterministically thinned in chronological order, whereas a secondary task uses at most 1,500 normal-domain source windows from each of the other five farms. The 128 labels are the model-fitting budget rather than the total target supervision. For Farms~1--6, the additional target-domain Q1 selection counts are 169, 617, 2,410, 212, 134, and 177, and the Q2 calibration counts are 305, 102, 862, 319, 362, and 225, respectively.}

{\color{blue}The formal-comparison input at each forecast contains four wind-speed measurements, sine and cosine encodings of four wind-direction measurements, temperature, atmospheric pressure, and relative humidity at each of the 24 strictly preceding 15-minute observations, giving 360 flattened meteorological inputs. Past power and future meteorological values are excluded. Power labels are divided by their farm-specific nominal capacity before fit-only label standardization and are mapped back to the target farm's MW scale for evaluation. Missing-value imputation and feature and label standardization are fitted only on the 2019 model-fitting rows. For pooled models, the source and target domains each contribute one half of the preprocessing weight; target-only preprocessing uses only its target fitting rows. Nystr\"om landmarks are constructed without labels from fitting inputs. A pooled landmark pool contains at most 2,000 source and 2,000 target rows and assigns total K-means weight one half to each domain; target-only construction uses only target fitting rows. K-means is used when the requested rank is smaller than the pool, while all pool rows are retained when the requested rank equals the pool size; any resulting effective-rank cap is recorded.}

{\color{blue}For each task--method pair, the frozen search rule issues 54 effective Q1 configuration queries through three predeclared sequential stages of 18 new queries: a base search, a middle active-regularizer search or method-equivalent refinement, and a local joint refinement. The middle stage either varies the active regularizer coefficients while retaining the base incumbent's other coordinates or applies a method-equivalent deterministic refinement around that incumbent; the joint stage is generated locally around the incumbent retained after the middle stage. Within a task--method pair, the 18 effective configurations in each stage are distinct from one another and from every earlier-stage query, giving 54 distinct effective queries across the three stages. Kernel candidates are scored by the mean target-Q1 RMSE of the $\tau=0.5$ fit across seeds 0, 1, and 2, whereas deterministic HistGBR is fitted once at seed 0 per query. The pooled kernel searches include requested Nystr\"om ranks 128, 256, and 512, and target-only SVQR uses ranks 32, 64, and 128. Inactive hyperparameter coordinates are ignored when effective-query uniqueness is checked. Equal query counts provide a common configuration-query budget; they do not imply equal model capacity, numbers of solver attempts, wall-clock cost, or exhaustive optimization.}

{\color{blue}Candidate generation, model fitting, preprocessing, landmark construction, and Q1 scoring are completed before the immutable selection lock is created, and the analysis specification is separately frozen before any formal-comparison Q3--Q4 output is read. Each locked configuration is refitted on the unchanged 2019 fitting rows; Q1 labels are not added to the fit. Q2 is reserved for interval calibration, and the formal workflow accesses Q3--Q4 only after the selection and analysis locks and then evaluates it under the locked settings. Kernel settings use landmark seeds 0, 1, and 2, and deterministic HistGBR uses seed 0. No task, method, seed, or configuration is selected from Q2 or Q3--Q4 performance. Because the public-data test period had already been examined in earlier work, this design provides test-locked evidence on a reused period rather than independent external confirmation.}

{\color{blue}Every kernel fit is first attempted for at most 3,000 ADMM iterations. The numerical certificate is checked every 25 iterations and requires $P-D\leq 10^{-7}+10^{-3}\max\{1,|P|,|D|\}$. If the first attempt misses this condition, a second attempt is warm-started from its returned best-primal $(\boldsymbol{w},b)$ incumbent, with freshly initialized split and dual states, and is allowed at most 12,000 additional ADMM iterations; the data, feature map, objective, and hyperparameters remain unchanged. The candidate remains in the audit record but is ineligible for selection if any required selection seed still misses the certificate. Under the original execution rule, a locked main-study or matched-ablation fit that still missed the certificate produced no retained result. This is a numerical primal--dual-gap certificate for the chosen fixed-feature objective, not an interval-arithmetic proof, a certificate of Nystr\"om approximation error, or a predictive-generalization guarantee.}

{\color{blue}One locked main-study fit exhausted these two attempts with a relative gap of 0.0010217, compared with the unchanged 0.001 tolerance. Before its Q3--Q4 predictive scores were inspected, we froze a performance-blind solver amendment that committed the hashes of all then-existing completion artifacts and the already frozen analysis specifications. For an otherwise exact locked main or matched fit with no completed result, the amendment permitted one further warm restart, from the same best-primal incumbent and fresh split and dual states, for at most 96,000 iterations under the same objective, data, feature map, hyperparameters, and certificate tolerance. No completed result could be replaced. The single triggering main-study fit then met the unchanged certificate; no matched-ablation fit required the amendment. Selection, calibration, prediction, and analysis rules were unchanged, and the original failed record, model, amendment lock, and completion sidecars are retained in the reproducibility package.}
% END SELECTION BLOCK selection_protocol

% BEGIN SELECTION BLOCK selection_results
\subsubsection{Locked point-prediction results}
{\color{blue}Table~\ref{tab:selection-v2-primary-within} reports Q3--Q4 RMSE on every primary within-site task, and Table~\ref{tab:selection-v2-scope} reports the site-equal capacity-normalized macro RMSE and average task rank separately for the primary and secondary scopes. RMSE is calculated separately for every predeclared evaluation seed and then averaged; predictions from different seeds are neither ensembled nor treated as independent observations. Figure~\ref{fig:selection-v2-within-rmse} displays the six primary tasks, and Fig.~\ref{fig:selection-v2-scope-rank} shows the scope-level average ranks.}

\input{selection_tables}

\begin{figure}[pos=htbp]
\centering
\includegraphics[width=\linewidth]{figures/within_site_normalized_rmse.pdf}
\caption{\textcolor{blue}{Capacity-normalized Q3--Q4 RMSE on the six primary within-site normal-to-extreme tasks. All methods and prespecified seeds are retained. For a kernel method, the marker is the mean of its three separate-seed RMSE values and the error bar spans their minimum and maximum; HistGBR is deterministic. These seed ranges describe algorithmic sensitivity and are not confidence intervals.}}
\label{fig:selection-v2-within-rmse}
\end{figure}

\begin{figure}[pos=htbp]
\centering
\includegraphics[width=\linewidth]{figures/scope_average_rank.pdf}
\caption{\textcolor{blue}{Average task rank in the primary within-site and secondary cross-site scopes. Ranks are computed from each task's seed-mean RMSE before equal averaging over the six sites in a scope; lower values are better.}}
\label{fig:selection-v2-scope-rank}
\end{figure}

% @@RESULTS_POINT@@
% Insert one verified paragraph naming the empirical macro-RMSE winner, TL-QLDMR's
% macro endpoint and average rank, and its task wins (including exact ties) in each scope. Describe the
% primary and secondary scopes separately and do not turn rank into a significance claim.
% END SELECTION BLOCK selection_results

% BEGIN SELECTION BLOCK selection_point_inference
\subsubsection{Paired temporal-block inference}
{\color{blue}For every task and timestamp, we first average squared error over each method's declared evaluation seeds and then subtract the comparator's value from the TL-QLDMR value; negative differences favor TL-QLDMR. The timestamps are partitioned into non-overlapping Monday--Sunday seven-calendar-day blocks. We use 10,000 whole-block bootstrap resamples for a percentile 95\% interval and 99,999 two-sided paired block sign flips with a plus-one correction. Holm adjustment is applied in two separate families of 30 tests, one for the six primary tasks and one for the six secondary tasks, with five TL-QLDMR comparisons per task. The bootstrap treats block effects as approximately independent and exchangeable, while the sign-flip test additionally requires approximate sign symmetry under the null.}

{\color{blue}As a complementary scope-level summary, all 15 method pairs are compared using 99,999 two-sided seven-day block swaps of the exact site-equal capacity-normalized seed-mean macro-RMSE endpoint, with Holm adjustment applied separately within each scope. When a pair contains deterministic HistGBR, its single timewise error curve is compared with each kernel seed curve under the same block swaps; this preserves the declared seed-mean endpoint and does not treat the repeated deterministic curve as independent replication. The conservative first tier contains the empirical endpoint winner and every method whose inferiority to that winner is not established by the predeclared strict adjusted-$p<0.05$ exclusion rule. Membership indicates non-rejection of inferiority under this block procedure; it is not evidence of equivalence, does not establish superiority, and is not a formal model-confidence set.}

% @@RESULTS_FIRST_TIER@@
% Insert the verified empirical winner and complete first-tier membership for each
% scope, including TL-QLDMR whether it is retained or excluded.

% @@RESULTS_TASK_INFERENCE@@
% Insert only task-level TL-QLDMR-versus-comparator findings supported by the completed
% separate 30-test Holm families. Retain null and adverse findings when material.

{\color{blue}Table~\ref{tab:selection-v2-full-inference} retains all task-level TL-QLDMR comparisons, including nonsignificant and adverse effects. The test period is reused and the number of temporal blocks is limited, so these paired results are interpreted conservatively.}

% END SELECTION BLOCK selection_point_inference

% BEGIN SELECTION BLOCK ablation
\subsubsection{Matched component removal}
{\color{blue}We use a frozen matched ablation on all twelve tasks to separate component removal from retuning. The four matched configurations, their shared fit-only artifact hashes, and the matched-analysis specification are locked before any main or matched Q3--Q4 evaluation. This point-only analysis uses $\tau=0.5$ and no Q2 labels. Within each task and seed, the TL-QLDMR-selected kernel bandwidth, total loss mass, target fraction, requested rank, fitting rows, preprocessing, landmarks, and feature map are held fixed. The no-variance, no-discrepancy, and weighted-SVQR variants differ only by setting $\lambda_1$, $\lambda_2$, or both coefficients to zero. No additional Q1 query is made. This design estimates removal effects conditional on the configuration selected for TL-QLDMR and is distinct from the independently tuned procedure comparison above.}

{\color{blue}The matched analysis retains all four variants, all twelve tasks, and seeds 0, 1, and 2. It reports site-equal capacity-normalized macro RMSE and mean rank among the four matched variants, separately for the two scopes. At the task level, we subtract each removal variant's seed-averaged squared error from that of TL-QLDMR; negative differences favor TL-QLDMR. The resulting differences use the same seven-calendar-day blocks, 10,000 whole-block bootstrap resamples, and 99,999 two-sided sign flips, with Holm adjustment in separate families of 18 tests for the primary and secondary scopes. Failure to reject a removal effect is not evidence of equivalence, and a conditional removal effect does not establish the ordering of the independently retuned procedures.}

\input{matched_ablation}

% @@RESULTS_MATCHED_ABLATION@@
% Insert the verified aggregate matched effects, ranks, and Holm-supported task-level
% findings. State conditional interpretation and retain adverse/null component effects.
% END SELECTION BLOCK ablation

% BEGIN SELECTION BLOCK selection_optimization_results
\subsubsection{Optimization accounting}
{\color{blue}The frozen design contains 72 task--method selections and 54 newly queried effective configurations per selection, for 3,888 Q1 configuration queries. The five kernel procedures require three selection-seed fits per query and HistGBR one deterministic fit, giving 10,368 required validation fits before any certificate retry. A carried incumbent may compete at a later stage using its cached Q1 score but is neither refitted nor counted as a new query. Appendix Table~\ref{tab:selection-v2-optimization-summary} accounts for every newly queried candidate and reports, by scope and method, the query and eligibility counts, failed seed certificates, solver attempts, iterations, and recorded fitting time. Because task workers can run concurrently and feature maps are cached across variants, summed fit time is descriptive and is not used to rank deployment speed.}

% @@RESULTS_OPTIMIZATION@@
% Insert only verified certificate totals and descriptive timing from the completed
% optimization-summary artifact. Do not infer approximation quality or deployment throughput.
% END SELECTION BLOCK selection_optimization_results

% BEGIN SELECTION BLOCK probabilistic_multiquantile
\subsubsection{Legacy exploratory multi-quantile analysis}
{\color{blue}The within-Farm~6 nine-quantile experiment is retained as a separate exploratory probability sensitivity analysis. It predates the frozen formal protocol and uses the earlier 13-variable, 24-lag input track, giving 312 meteorological inputs. In this track, pooled methods use 7,500 source fitting windows, and every method uses all 2,076 target fitting labels. The archived legacy artifacts label TL-QLDMR as ``Full''. The predefined quantile grid is}
\begin{equation*}
{\color{blue}\mathcal{T}=\{0.025,0.05,0.10,0.25,0.50,0.75,0.90,0.95,0.975\}.}
\end{equation*}

{\color{blue}Each kernel procedure receives one reference configuration $(\gamma,A,p,\lambda_1,\lambda_2)=(0.003,1000,0.5,1,1)$ and 15 fixed Latin-hypercube configurations. Their ranges are $\gamma\in[10^{-4},10^{-1}]$, total loss mass $A\in[10^2,10^5]$, target fraction $p\in[0.05,0.95]$, and active regularizer coefficients in $[10^{-2},10^3]$; all coordinates except $p$ use logarithmic scales. The Latin-hypercube design uses fixed seed 20261022, and every legacy kernel candidate uses 128 requested landmarks. Source and target loss coefficients are $C_S=A(1-p)/n_S$ and $C_T=Ap/n_T$, inactive regularizer coordinates are ignored, and target-only SVQR uses $p=1$ and $C_T=A/n_T$. With the same fixed seed, HistGBR receives 16 configurations sampled without replacement from the Cartesian grid formed by learning rates $\{0.03,0.1\}$, leaf limits $\{7,15,31\}$, minimum leaf counts $\{10,20,40\}$, $L_2$ weights $\{0,1,10\}$, iteration counts $\{100,200\}$, and target fractions $\{0.1,0.5,0.9\}$. This legacy track uses fitting-only equal-row preprocessing, so domain influence follows the available row counts, and ordinary unweighted K-means on the sampled landmark pool; it does not use the explicit one-half source/one-half target weights of the formal comparison.}

{\color{blue}Candidates are scored by unweighted mean Q1 pinball loss over all nine quantiles. Ties use the fixed candidate index, and the archived legacy rule does not filter candidates by convergence status. The 96 method--candidate records comprise 864 seed-0 quantile fits, and all six settings are locked before any Q2 calibration or Q3--Q4 prediction. The selected kernel settings are then evaluated at landmark seeds 0, 1, and 2, while deterministic HistGBR uses seed 0. Equal candidate counts specify this legacy probability-selection budget but do not imply equal model capacity or computational cost. The six selected settings are reported with the legacy audit material in Appendix Table~\ref{tab:probabilistic_configs}.}

{\color{blue}Raw quantile forecasts are retained for quantile-specific pinball loss, empirical reliability, and crossing diagnostics. A fixed label-free monotone rearrangement sorts the nine forecasts within each row before nested intervals and the weighted interval score are calculated. Q2 scores provide a fixed split-conformal expansion separately for the 50\%, 80\%, 90\%, and 95\% central intervals. The past-only rolling rule starts from Q2, uses at most the 256 most recent scores available before each forecast, and appends the current Q3--Q4 score only after that forecast has been issued. No Q3--Q4 label selects a configuration or the monotone-rearrangement rule; the fixed expansion uses Q2 alone, whereas each rolling expansion may use outcomes revealed strictly before that forecast.}

{\color{blue}The archived selection lock records the six selected configurations, and the saved validation, calibration, and test artifacts are covered by a post-completion integrity inventory. The original run manifest did not freeze the temporal loader, the target-sampling helper, the source-workbook hashes, the chronological-reference manifest, or the complete software environment; it recorded only selected source hashes and the Python and NumPy versions. This legacy analysis also uses the previously examined Q3--Q4 period. We therefore classify it as legacy/post hoc exploratory evidence rather than part of the formal point comparison or matched component analysis; its pinball loss, WIS, reliability, crossing, coverage, width, and paired-WIS results do not provide independent confirmation.}

{\color{blue}Table~\ref{tab:probabilistic_scores} shows that HistGBR has the lowest raw nine-quantile mean pinball loss (1.767\,MW) and rearranged WIS (3.531\,MW), followed by target-only SVQR (1.995 and 3.971\,MW). TL-QLDMR (Full) gives 2.045 and 4.090\,MW. Its raw crossing rate is zero, as is that of the no-discrepancy variant, whereas the other four methods cross on 38.07\%--45.84\% of test rows. This ordering stability does not translate into the lowest proper score or calibration error: TL-QLDMR's rearranged reliability MACE is 8.66 percentage points, compared with 7.35 for the no-variance and weighted-SVQR fits. These results do not establish a general probabilistic advantage for the complete objective.}

{\color{blue}For TL-QLDMR (Full), fixed Q2 calibration gives PICP values 0.532, 0.894, 0.960, and 0.977 at nominal coverages 0.50, 0.80, 0.90, and 0.95, with widths 11.21, 30.86, 41.52, and 54.75\,MW (Table~\ref{tab:probabilistic_coverage}). Figure~\ref{fig:probabilistic_reliability} reports all nine reliability points before and after rearrangement, while Fig.~\ref{fig:probabilistic_coverage} shows the static coverage--width trade-off without selecting methods or levels by test performance. Past-only rolling results are retained in the table as a separate information setting.}

{\color{blue}Paired probability-score inference uses the same 25 seven-day calendar blocks. WIS is computed separately for every seed and forecast, then averaged across landmark seeds within that forecast before forming paired method differences. The WIS difference (TL-QLDMR minus comparator) is $-0.0009$\,MW against the no-discrepancy variant (95\% block-bootstrap interval $[-0.0014,-0.0004]$, Holm $p=0.0058$), a statistically detectable but practically negligible difference. Against HistGBR the difference is $0.5585$\,MW ($[0.3290,0.8272]$, Holm $p=0.0020$), favoring HistGBR. The other three adjusted comparisons are not significant. Table~\ref{tab:probabilistic_inference} retains all five effects and tests. The block bootstrap treats block effects as approximately independent and exchangeable; the sign-flip test additionally requires approximate sign symmetry under the null. The small number of blocks and reused test period therefore require conservative interpretation.}


\input{probabilistic_scores}
\input{probabilistic_coverage}

\begin{figure}[pos=htbp]
\centering
\includegraphics[width=\linewidth]{figures/probabilistic_reliability.pdf}
\caption{\textcolor{blue}{Legacy Farm~6 reliability for the probability-selected configurations. All six methods and nine predefined levels are retained; panel (a) uses the independently fitted quantiles and panel (b) applies the fixed row-wise monotone rearrangement.}}
\label{fig:probabilistic_reliability}
\end{figure}

\begin{figure}[pos=htbp]
\centering
\includegraphics[width=\linewidth]{figures/probabilistic_coverage_sharpness.pdf}
\caption{\textcolor{blue}{Legacy Farm~6 fixed-Q2 interval calibration: empirical coverage and mean interval width at the four predefined nominal levels. Width is interpreted together with coverage.}}
\label{fig:probabilistic_coverage}
\end{figure}

\input{probabilistic_inference}
% END SELECTION BLOCK probabilistic_multiquantile

% BEGIN SELECTION BLOCK sequential_calibration
\subsubsection{Earlier reused-test interval calibration}
{\color{blue}In the earlier point-selected two-quantile diagnostic, two calibration protocols distinguish batch forecasting from sequential operation. The fixed protocol uses all Q2 target scores and the common expansion rule defined above. Past-only rolling calibration starts with the most recent 256 available Q2 scores (or all scores when fewer are available). Before target forecast $t$, the expansion is the $\lceil(n_t+1)0.95\rceil$-th ordered score in this buffer. Only after producing the interval is its observed outcome used to append a new score and discard the oldest score if needed. The point and quantile models remain fixed. All six methods use this same rule. The past-only rolling evaluation assumes earlier target outcomes are available before the next evaluated forecast and therefore consumes additional labels during testing; it is a different information setting from fixed calibration. Neither temporal protocol inherits an exchangeability-based coverage guarantee.}
% END SELECTION BLOCK sequential_calibration

% BEGIN SELECTION BLOCK selection_calibration_results
\subsubsection{Earlier exploratory cross-farm calibration}
{\color{blue}In this earlier point-selected diagnostic, fixed Q2 calibration for TL-QLDMR (Full in the archived artifacts) produces cross-farm PICP from 0.640 to 0.988. The lowest fixed-Q2 coverage occurs on Farm~5. Past-only recalibration raises its mean PICP from 0.640 to 0.917 and reduces mean Winkler score from 44.591 to 36.547\,MW, but it does not reach 0.95. Across the other five farms, rolling PICP ranges from 0.948 to 0.962. The adjustment can widen or narrow intervals as recent scores change: for example, the TL-QLDMR mean width on Farm~5 increases from 30.760 to 33.319\,MW, while that on Farm~6 decreases from 84.673 to 83.233\,MW. Figure~\ref{fig:selected_calibration} displays both coverage and capacity-normalized width. This exploratory comparison reports performance under past-only score updating; it is not a coverage guarantee, does not identify temporal drift as a cause, and uses more outcome information than fixed Q2 calibration.}

\begin{figure}[pos=htbp]
\centering
\includegraphics[width=\linewidth]{figures/selection_calibration.pdf}
\caption{\textcolor{blue}{Earlier exploratory TL-QLDMR (Full) fixed-Q2 and past-only rolling calibration across six target farms: (a) PICP and (b) mean interval width normalized by farm capacity. Points are three-seed means with minimum--maximum bars. The dashed line is nominal 0.95 coverage.}}
\label{fig:selected_calibration}
\end{figure}
\FloatBarrier
% END SELECTION BLOCK selection_calibration_results

\subsubsection{Earlier exploratory residual diagnostics}
\input{residual_diagnostics}

\FloatBarrier
